\subsection{群作逆紧作用的性质}

命题 3. 若拓扑群 G 逆紧作用于拓扑空间 X，则轨道空间 X/G 是豪斯多夫的。若 G 还是豪斯多夫的，则 X 也是豪斯多夫的。

设 \(C \subset X \times X\) 是等价关系 \(R\)（由 \(G\) 定义在 \(X\) 上的）的图；则 \(C\) 是 \(G \times X\) 在映射 \(\theta : (s, x) \to (x, s.x)\) 下的像。由于 \(\theta\) 是逆紧的，\(C\) 在 \(X \times X\) 中闭（第一章，§ 10，第 1 节，命题 1）。由于关系 \(R\) 是开的（§ 2，第 4 节，引理 2），由此推出 \(X/G\) 是豪斯多夫的（第一章，§ 8，第 3 节，命题 8）。

现在设 G 是豪斯多夫的。则映射 \(x \to (e, x)\)（X 到 \(G \times X\) 中的映射）是 X 到 \(G \times X\) 的一个闭子空间上的同胚，因而是逆紧的（第一章，§ 10，第 1 节，命题 2）。把这个映射与映射 \((s, x) \to (x, s.x)\)（\(G \times X\) 到 \(X \times X\) 中的映射，由假设它是逆紧的）复合，就得到 X 到 \(X \times X\) 中的一个逆紧映射，即对角映射 \(x \to (x, x)\)。于是对角线 \(\Delta\)（\(X \times X\) 的对角线）在 X 中是闭的，因此 X 是豪斯多夫的。

命题 4. 设 G 是逆紧作用于拓扑空间 X 的一个拓扑群，并设 x 是 X 的一点。用 G.x 表示 x 的轨道，用 \(K_{x}\) 表示 x 的稳定子群。那么：

\begin{enumerate}
\def\labelenumi{\alph{enumi})}
\item
  映射 \(s \to s.x\) 是 \(G\) 到 \(X\) 中的逆紧映射。
\item
  \(\mathbf{K}_x\) 是拟紧的。
\item
  \(\mathbf{G} / \mathbf{K}_x\) 到 \(\mathbf{G}.x\) 上的典范映射是同胚。
\item
  轨道 G.x 在 X 中是闭的。
\end{enumerate}

\(\{x\} \times X\) 在 \(\theta : (s, x) \to (x, s.x)\) 下的逆像是 \(G \times \{x\}\)；因此，由第一章 § 10 第 1 节命题 3，\(\theta\) 在 \(G \times \{x\}\) 上的限制是 \(G \times \{x\}\) 到 \(\{x\} \times X\) 中的逆紧映射，由此得 a)。由于 \(K_x\) 是 \(x\) 在 \(s \to s.x\) 下的逆像，故 b) 由第一章 § 10 第 2 节定理 1 得出。c) 与 d) 是 a) 的推论，这依据第一章 § 10 第 1 节命题 2 与命题 5 b)。

注. 命题 4 表明，若拓扑群 G 逆紧作用于齐性空间 G/H，则子群 H 是拟紧的（从而当 G 是豪斯多夫的时它是紧的）。可以证明，这也是 G 逆紧作用于 G/H 的一个充分条件（习题 3）。

命题 5. 设 G（相应地，G'）是连续作用于拓扑空间 X（相应地，X'）的拓扑群。设 φ 是 G 到 \(\mathbf{G}'\) 中的一个连续同态，并设 \(\psi\) 是 X 到 \(\mathbf{X}'\) 中与 \(\varphi\) 相容的连续映射（ \(\S 2\) ，第 4 节）。那么：

\begin{enumerate}
\def\labelenumi{(\roman{enumi})}
\item
  若 \(\varphi\) 是满射，\(\psi\) 是满射且逆紧的，并且 \(\mathbf{G}\) 逆紧作用于 \(\mathbf{X}\)，则 \(\mathbf{G}'\) 逆紧作用于 \(\mathbf{X}'\)。
\item
  若 \(\varphi\) 是逆紧的，\(\mathbf{G}'\) 逆紧作用于 \(\mathbf{X}'\)，且 \(\mathbf{X}\) 是豪斯多夫的，则 \(\mathbf{G}\) 逆紧作用于 \(\mathbf{X}\)。
\end{enumerate}

为证 (i)，考虑交换图

\[
\begin{array}{c c c c} \mathbf {G} & \times \mathbf {X} \xrightarrow {\theta} \mathbf {X} & \times \mathbf {X} \\ \alpha & \downarrow & & \downarrow \beta \\ \mathbf {G} ^ {\prime} & \times \mathbf {X} ^ {\prime} \xrightarrow {\theta^ {\prime}} \mathbf {X} ^ {\prime} & \times \mathbf {X} ^ {\prime} \end{array}
\]

其中 \(\alpha = \varphi \times \psi\)，\(\beta = \psi \times \psi\)。由假设，\(\theta\) 是逆紧的；\(\beta\) 也是逆紧的{[}第一章，§ 10，第 1 节，命题 4a){]}；于是 \(\beta \circ \theta = \theta' \circ \alpha\) 是逆紧的{[}第一章，§ 10，第 1 节，命题 5a){]}。由于 \(\alpha\) 是满射，由此推出 \(\theta'\) 是逆紧的{[}第一章，§ 10，第 1 节，命题 5b){]}。

为证 (ii)，考虑超滤子 \(\mathfrak{U}\)，它定义在 \(\mathbf{G} \times \mathbf{X}\) 上，使得映射

\[
(s, x) \rightarrow s. x \quad \text { and } \quad (s, x) \rightarrow x
\]

关于 \(\mathfrak{U}\) 分别收敛于 \(y_0\) 与 \(x_0\)。由此推出 \((s, x) \to \varphi(s). \psi(x)\) 与 \((s, x) \to \psi(x)\) 关于 \(\mathfrak{U}\) 收敛。由于 \(\mathbf{G}'\) 逆紧作用于 \(\mathbf{X}'\)，这就意味着（第 1 节）\((s, x) \to \varphi(s)\) 关于 \(\mathfrak{U}\) 收敛于一点 \(s_0' \in \mathbf{G}'\)。由于 \(\varphi\) 是逆紧的，我们推出（第一章，§ 10，第 2 节，定理 1）\((s, x) \to s\) 关于 \(\mathfrak{U}\) 收敛于一点 \(s_0 \in \mathbf{G}\)。于是在 \(\mathbf{X}\) 中极限的唯一性表明 \(y_0 = s_0 x_0\)，从而 \(\mathbf{G}\) 逆紧作用于 \(\mathbf{X}\)（第 1 节）。

